% =====================================================================
% journey_detailed.tex - THE COMPLETE JOURNEY, checkpoint by checkpoint.
% The most detailed document in the pack: a reader who reads only this
% should understand the work AND everything faced along the way.
% Style: first person, simple language, every jargon term explained at
% first use. Compile: pdflatex journey_detailed.tex (twice for TOC).
% =====================================================================
\documentclass[11pt]{article}
\usepackage[a4paper,margin=1in]{geometry}
\usepackage{amsmath,amssymb}
\usepackage{booktabs}
\usepackage{enumitem}
\usepackage{framed}
\usepackage[hidelinks]{hyperref}
\setlength{\parskip}{4pt}

\newenvironment{checkpoint}[1]
  {\begin{framed}\noindent\textbf{CHECKPOINT - #1.}\ }
  {\end{framed}}
\newenvironment{lesson}
  {\begin{framed}\noindent\textbf{What this taught me.}\ \itshape}
  {\end{framed}}

\title{\textbf{The Complete Journey}\\[4pt]
\large Evaluating India's 2021 Agricultural Derivatives Ban -\\
every checkpoint, every error, every fix, from first question to final number}
\author{}\date{}

\begin{document}
\maketitle
\tableofcontents
\newpage

% =====================================================================
\section{How to read this document}

This is the long version of my project story - the one that contains everything. It is
organised the way the work actually happened, in four parts:

\begin{enumerate}[leftmargin=1.8em,itemsep=1pt]
  \item \textbf{The question and the rules} - what I set out to test, and the discipline I
        fixed before touching data.
  \item \textbf{Data collection} - honestly about half the project. How I got 5.7 million
        price records out of systems that were never designed to hand them over: the API
        work, the automation, the network errors, the traps.
  \item \textbf{Data cleaning} - how raw downloads became an analysis-ready panel, with
        the actual Python functions that did it.
  \item \textbf{Execution} - the estimation arc: difference-in-differences, the event
        study, how my design failed its own tests, the data error that failure exposed,
        the fix, the honest inference, and the comparison-group revision that changed my
        answer.
\end{enumerate}

I use technical terms freely but explain each one in plain words the first time it
appears. That is deliberate: being able to say things simply is part of knowing them.

% =====================================================================
\section{Part 1 - The question and the rules}

\subsection{The policy and the question}

In December 2021, India's market regulator SEBI suspended futures trading in seven
agricultural commodities - wheat, paddy (rice), chana (chickpea), mustard, soybean, moong,
and crude palm oil. A \textbf{futures contract} is an agreement to buy or sell a commodity
at a fixed price on a future date; farmers and traders use futures to lock in prices, and
the futures price doubles as everyone's best public forecast of where the market is going.
The government's premise was that speculation in these contracts was inflating actual
market prices - so switching them off should calm things down.

That premise is testable. The ban is a sharp, dated event that hit some commodities and
not others, which sets up a natural comparison:

\begin{checkpoint}{The research question}
Did spot-price volatility - the day-to-day jumpiness of real wholesale (\textbf{mandi})
prices - fall for the banned commodities after the ban, relative to comparable commodities
whose futures kept trading? And whatever the answer, what did the ban cost in market
function?
\end{checkpoint}

Two vocabulary items used throughout. The \textbf{spot price} is the price for immediate
delivery - what the commodity actually sells for in a mandi today (as opposed to the
futures price, which is for later delivery). \textbf{Volatility} here means realized
volatility: take daily percentage price changes, compute their standard deviation over a
rolling 30-day window, and scale it to an annual figure. High volatility = prices jumping
around; low = stable.

\subsection{The rules I fixed before starting}

Empirical work is easy to fool yourself in, so I locked five rules up front:

\begin{enumerate}[leftmargin=1.8em,itemsep=1pt]
  \item \textbf{Raw data is immutable.} Every downloaded file is stored exactly as
        received and never edited. All changes happen in scripts, so any transformation
        can be inspected and re-run.
  \item \textbf{Every acquisition is logged.} One row per data series in a provenance log:
        source, URL, dates, units, known issues.
  \item \textbf{Pre-registration.} I wrote down the volatility measure, the estimating
        equation, and explicit pass/fail rules - including falsification tests the design
        must survive - \emph{before} estimating anything. (Why this matters becomes the
        centre of the story in Part 4.)
  \item \textbf{Hard truncation of banned futures} at 20 December 2021 - for a reason I
        only fully appreciated once I saw the raw exchange data (Part 2.6).
  \item \textbf{Dated decisions.} Every methodological choice got written down with its
        rationale when it was made, so the path from first idea to final specification is
        auditable - including the choices I later had to revise.
\end{enumerate}

% =====================================================================
\section{Part 2 - Data collection (half the project)}

There was no ready-made dataset for this question. Every series had to be extracted from
a system that was not designed to hand it over in bulk. This part is long because the
work was long.

\subsection{Finding a route to mandi prices}

India's official wholesale price portal, \textbf{Agmarknet}, publishes prices for every
mandi (regulated wholesale market) - but through web forms designed for a human looking up
one commodity, one market, one date range at a time. I needed nine years, twenty
commodities, hundreds of districts. Clicking was not an option, and scraping thousands of
form submissions is slow and fragile.

The workable route: the Centre for Economic Data and Analysis (\textbf{CEDA}) at Ashoka
University re-serves the same Ministry of Agriculture data, cleaned, through a documented
\textbf{API}. An API - Application Programming Interface - is a way for a program to
request data from a server directly: instead of a human loading a webpage, my script sends
a structured request to a fixed address (an \textbf{endpoint}) and the server replies with
structured data. Registration is free for academic use; it issues an \textbf{API key} - a
long secret string that identifies my account, which my script presents with every
request. For provenance I cite both CEDA and the underlying Agmarknet source.

\subsection{How the API conversation actually works}

Each request is an \textbf{HTTP POST}: my script sends a small \textbf{JSON} package -
JavaScript Object Notation, the standard text format for structured data, all curly braces
and key-value pairs - to the prices endpoint. The request body says what I want:

\begin{verbatim}
POST https://<host>/v1/agmarknet/prices
Authorization: Bearer <my API key>          <- the key rides in a request header
{ "commodity_id": 6,                        <- chana
  "state_id": 0,                            <- 0 = national
  "from_date": "2017-01-01",
  "to_date":   "2026-04-30",
  "indicator": "price" }
\end{verbatim}

A \textbf{header} is metadata attached to the request; the \texttt{Authorization: Bearer}
header is the standard place an API key travels. The server replies with a JSON payload -
an array of rows (date, market, price) - plus its own \textbf{response headers} and a
\textbf{status code}: a three-digit summary of how the request went. The ones that shaped
my life for a week: \texttt{200} (success), \texttt{429} (too many requests - slow down),
\texttt{403} (forbidden - you look like a bot), \texttt{504} (gateway timeout - the server
gave up building your answer).

\subsection{Checkpoint: the rate-limit war}

\begin{checkpoint}{The 40-requests-per-hour wall}
The API allows only 40 requests per hour. My first downloader did not know that: it fired
requests, got \texttt{429} responses, retried blindly - and every retry \emph{also} counted
against the quota. It ran 17 minutes and produced zero rows, having silently burned the
entire hour's allowance on retries.
\end{checkpoint}

The fix is a small lesson in reading what the network tells you. Rate-limited APIs
advertise their state in response headers: \texttt{RateLimit-Remaining} (how many requests
you have left this window), \texttt{RateLimit-Reset} (seconds until the window renews), and
on a \texttt{429}, \texttt{Retry-After} (how long to wait). I rebuilt the downloader to be
\textbf{rate-limit aware}: after every response it reads these headers, and when the
remaining count hits zero it sleeps exactly until the reset instead of hammering the
server. Polite, and optimal - you extract precisely the allowed 40/hour, no more requests
wasted on rejections.

\subsection{Checkpoint: architecture for a ten-hour download}

At 40 requests an hour, the full district-level pull - every commodity, every state -
needs hundreds of requests, meaning \textbf{roughly ten hours} of wall-clock time. Any
crash, network blip, or laptop sleep must not restart the job from zero. So the
downloader's architecture:

\begin{itemize}[leftmargin=1.8em,itemsep=1pt]
  \item \textbf{One JSON file per unit of work} - one commodity x one state - named
        predictably (\texttt{chana\_\_state8.json}). Before requesting anything, the
        script checks whether that file already exists; if yes, skip. That single rule
        makes the whole job \textbf{resumable}: kill it anywhere, relaunch, and it
        continues where it stopped.
  \item \textbf{Self-describing files.} Every saved JSON carries a \texttt{\_meta} block -
        source, commodity id, state, date range, pull date, row count - followed by the
        raw data array. Months later, any file explains itself without consulting notes.
  \item \textbf{Transient-error tolerance.} Real networks hiccup: connections reset,
        sockets time out. The script retries these with escalating pauses - 8 seconds,
        then 16, then 24, capped at one minute (\textbf{backoff}; the textbook variant
        that doubles each wait is called exponential backoff) - and only gives up on a
        request after eight consecutive failures, rather than dying at 3 a.m. on hour
        seven.
\end{itemize}

The pull completed across all 396 commodity-state files: \textbf{5.74 million market-day
rows}, zero dropped. The design also passed a live-fire test I had not planned: the
terminal crashed mid-run one night. Relaunching cost nothing - every finished file was
already on disk, the script skipped them and carried on - and afterwards I armed an
hourly scheduled health-check (a \textbf{cron} job: the standard Unix timer that runs a
command on a schedule) to restart the download if it ever died unattended again.

\subsection{Checkpoint: the 504 problem, solved by divide and conquer}

\begin{checkpoint}{Big states broke the server}
For large states like Uttar Pradesh, asking for nine years x all districts in one request
made the server time out: \texttt{504 Gateway Timeout} - it could not assemble that much
data before its own internal deadline.
\end{checkpoint}

The response payload was simply too big, so the fix is to ask for less at a time, adaptively:

\begin{enumerate}[leftmargin=1.8em,itemsep=1pt]
  \item Request districts in \textbf{batches} (about 15 at a time) instead of a whole state.
  \item If a batch still times out, \textbf{halve it} and try each half - recursively.
  \item If even a \emph{single district} times out (Raipur's paddy data was so dense the
        district overwhelmed the server on its own), switch dimension: \textbf{halve the
        date range} instead, recursively, until the server can serve each piece.
\end{enumerate}

This is divide-and-conquer applied to a network problem: keep splitting the request along
whichever dimension is too big until every piece fits under the server's timeout. No data
is abandoned; it just arrives in smaller boxes.

\subsection{Checkpoint: three traps in the data itself}

\paragraph{Trap 1 - identifier mismatch.} The API's commodity identifiers do \emph{not}
match the official portal's identifiers for the same crops (turmeric is 39 in one system
and 35 in the other; jeera 42 vs 38). If I had assumed they matched, I would have
downloaded the wrong commodities with full confidence. Fix: query the API's own
\texttt{/commodities} listing first and verify every identifier by name before pulling.

\paragraph{Trap 2 - the exchange that hates robots.} Crude palm oil trades on the MCX
exchange, whose website returns \texttt{403 Forbidden} to every non-browser client -
scripts, command-line tools, everything. Sites detect this mainly through the
\textbf{User-Agent} header (the string identifying what software is asking) and other
browser fingerprints. The workaround: drive a \emph{real} browser session and use the
exchange's own data endpoint - the page's JavaScript revealed exactly which internal
address it calls for price tables, and calling that from within a genuine browser context
is indistinguishable from the website's own behaviour. That yielded \textbf{all 64
palm-oil contracts spanning 2017 to the suspension}, downloaded with zero failures, each
with traded volume and open interest - fields no vendor provides.

\paragraph{Trap 3 - the ghost prices.} Inside that raw exchange data sat the project's
most important warning: contracts kept publishing daily ``settlement'' marks for
\emph{months after the ban} - rows with prices but zero volume, zero trades. They look
like post-ban futures prices; they are administrative echoes. This is precisely why rule 4
exists: every banned-commodity futures series is hard-truncated at 20 December 2021 in the
cleaned data, so no ghost quote can ever masquerade as a market price.

\subsection{Checkpoint: the guar forensics - validate everything against a benchmark}

\begin{checkpoint}{A mislabelled series, caught by correlation}
I cross-checked every spot series against its own futures series where one existed - two
views of the same commodity should co-move. Every commodity passed (correlations
0.92-0.99) except guar: \textbf{0.04}. Essentially uncorrelated with its own futures.
\end{checkpoint}

Investigation: the series labelled ``Guar'' (identifier 75) silently mixed \textbf{guar
gum} - a processed derivative roughly \emph{ten times} the price of the raw seed - in with
guar seed. The blend was meaningless as a price series. The database held a separate,
correctly-scoped ``Guar Seed'' series (identifier 413); its correlation with guar futures:
\textbf{0.99}. I switched, and documented both series so the audit trail shows the swap.

\begin{lesson}
Never trust a label. Every series earns its place by agreeing with an independent
benchmark. This one check later justified itself many times over.
\end{lesson}

\subsection{The other sources, briefly}

\begin{itemize}[leftmargin=1.8em,itemsep=1pt]
  \item \textbf{Vendor futures for comparison commodities} (castor, guar, jeera, turmeric,
        kapas): the exchange's own public archive turned out to reach back only to July
        2024 - the first wall the project hit - so these came from a financial-data vendor
        instead: daily prices back to 2017, pulled as three \textbf{contract legs} - c1
        (front month), c2 (next), c3 (far) - because liquidity differs along the curve.
        A continuous futures series is stitched from expiring contracts, and \emph{how}
        the vendor stitches matters; I \textbf{reverse-engineered their rule from the data}
        (detect days where the front series suddenly adopts the next contract's previous
        price - those are roll days). Verdict: hold the front contract through expiry,
        switch the next session, no adjustment at the splice. Knowing this let me build my
        own series identically, so constructions are comparable.
  \item \textbf{SEBI monthly bulletins} (who trades these markets): the listing pages are
        JavaScript shells with no scrapable links, but the document pages themselves are
        served to standard web crawlers - so the document URLs were recoverable via search
        engines and web archives, then downloaded from a frozen list. 33 of 36 months
        (three are missing on the regulator's own server).
  \item \textbf{International references}: CBOT wheat futures and the dollar-ringgit
        exchange rate, via a public financial-data interface.
  \item \textbf{The policy ledger}: 24 dated interventions (the ban, its extensions,
        export bans, stock limits, duty changes), 22 of the 24 pinned to a primary
        government document with the URL recorded. This matters enormously later: the ban
        did not happen in isolation, and this ledger is how I keep other policies from
        being mistaken for the ban's effect.
  \item \textbf{The food-donor extension} (added mid-project for reasons Part 4 explains):
        barley, maize, jowar pulled at national and district level (bajra's district pass
        nearly complete); groundnut, sesamum, sunflower, ragi at national level with
        district downloads still queued behind the rate limit.
\end{itemize}

% =====================================================================
\section{Part 3 - Data cleaning: from raw JSON to an analysis panel}

\subsection{The philosophy and the three tiers}

Nothing is ever edited by hand. Data lives in three tiers: \texttt{raw/} (verbatim
downloads), \texttt{constructed/} (deterministic derivations, like the stitched futures
series), and \texttt{clean/} (what estimation reads) - and only scripts move data between
tiers. Every cleaning script prints diagnostics on every run - date span, rows per year,
largest gap, count of extreme daily moves, price percentiles - so anomalies surface
immediately rather than hiding inside a CSV.

\subsection{The cleaning rules, and the code that enforces them}

The heart of the pipeline is a handful of small functions - small enough to show.

\paragraph{Parsing the API's JSON.} Step zero is turning each saved JSON file back into
a table. Verbatim from the panel builder:
\begin{verbatim}
blob = json.loads(fp.read_text())     # {'_meta': {...}, 'data': [rows]}
df = pd.DataFrame(blob["data"])       # list of row-dicts -> DataFrame
df["date"] = pd.to_datetime(df["date"].str[:10])
df = df[(df.modal_price > 0) & trading_days_only(df.date)]
\end{verbatim}
Three cheap guards stand at the door: dates parsed from the first ten characters of the
API's timestamp string, non-positive prices dropped, and the weekday filter (whose story
is below) applied before anything is computed.

\paragraph{Log returns.} Daily price changes are computed as differences of logarithms
(log returns), the standard because they are symmetric and additive over time:
\begin{verbatim}
def log_returns(price):
    return np.log(price).diff()      # ln(P_t) - ln(P_{t-1})
\end{verbatim}

\paragraph{Realized volatility.} A rolling 30-day standard deviation of those returns,
scaled to annual units by $\sqrt{252}$ - because there are about 252 trading days in a
year:
\begin{verbatim}
def realized_vol(returns, window=30, annualize=True):
    rv = returns.rolling(window, min_periods=int(window*0.8)).std()
    return rv * np.sqrt(252) if annualize else rv
\end{verbatim}
(\texttt{min\_periods} tolerates a few missing days without producing fake gaps.)

\paragraph{Trading days only.} Keep weekday observations before computing anything:
\begin{verbatim}
def trading_days_only(dates):
    return pd.to_datetime(dates).dt.dayofweek < 5   # Mon-Fri
\end{verbatim}
This one-line filter exists because of the most important error in the project - the full
story is in Part 4. I state the rule here because it is now part of cleaning; I found the
\emph{reason} for it the hard way.

\paragraph{Ban truncation.} The cleaner decides a commodity's ban status from its base
name, so leg-suffixed files (\texttt{cpo\_c1}) truncate exactly like their parent:
\begin{verbatim}
if commodity.lower().split('_')[0] in BANNED:
    df = df[df.date <= BAN_DATE]     # 2021-12-20; ghost quotes die here
\end{verbatim}

\paragraph{Robust aggregation.} District-level prices are built as \emph{medians}, not
means, at two levels: first the median across markets within a district-day, then (for the
national series) the median across districts. Why medians: a single fat-fingered mandi
entry - and I found one price recorded around 14 \emph{million} rupees against a median of
about 4{,}000 - destroys a mean but barely dents a median. The core pandas pattern:
\begin{verbatim}
dd = df.groupby(["census_district_id","date"])["modal_price"].median()
ret = panel.groupby(["state","district"])["price"].transform(log_returns)
monthly = (panel.set_index("date")
                .groupby(["commodity","state","district"])
                .resample("ME")[["rv30"]].mean())      # ME = month-end
\end{verbatim}
\texttt{groupby().transform()} applies a function \emph{within} each district's own series
(so returns never straddle two districts), and \texttt{resample("ME")} collapses daily
volatility to a monthly panel.

\paragraph{Price-support censoring.} For paddy, \textbf{40.3\%} of daily returns are
\emph{exactly zero} - because government procurement at the Minimum Support Price
(\textbf{MSP}, a floor price at which the state buys) pins the market price for weeks at a
time. ``Volatility'' of an administered price is a measurement artifact, so paddy is
excluded from the primary analysis; wheat (19.7\% zero returns) stays but carries a
censoring flag, and results are shown with and without it.

\subsection{The finished panel}

The output of all of this is one file: a commodity x district x month panel of realized
volatility - \textbf{172{,}381 cells across 16 commodities}, 2017 through October 2025,
with a \texttt{post} indicator switching on at the ban date. Every number in Part 4 is
computed from this panel, and the audit on it is clean: zero weekend rows, zero duplicate
dates, zero non-positive prices, uniform coverage.

% =====================================================================
\section{Part 4 - Execution: the estimation arc}

\subsection{The design: difference-in-differences, in plain words}

The naive approach - compare banned commodities' volatility before vs after - fails
because \emph{everything} changed in 2022 (a war, inflation, weather). The fix is a
control group. \textbf{Difference-in-differences (DiD)} takes two differences: the
before-to-after change for banned commodities, \emph{minus} the before-to-after change for
comparison commodities whose futures kept trading. Anything that hit both groups alike -
the war, the monsoon, fuel costs - subtracts out. Formally:
\[
\ln(\mathrm{vol}_{it}) = \alpha_i + \lambda_t + \beta\,(\mathrm{banned}_i \times
\mathrm{post}_t) + \varepsilon_{it}
\]
where $\alpha_i$ are \textbf{unit fixed effects} (one constant per commodity-district,
absorbing permanent differences - turmeric is just always jumpier than wheat),
$\lambda_t$ are \textbf{month fixed effects} (absorbing anything common to all commodities
in a month), and $\beta$ is the answer. Because the outcome is in logs, $\beta$ converts
to a percentage: effect $= e^{\beta}-1$.

One inference decision mattered enormously: the ban was assigned at the \emph{commodity}
level, so however many district rows I have, there are really only about a dozen
independent policy units. Standard errors are therefore \textbf{clustered by commodity} -
meaning the uncertainty calculation treats each commodity as one correlated block rather
than pretending thousands of district-months are independent evidence.

\subsection{The falsification battery - designed before estimation, on purpose}

Before running anything, I committed to two tests the design had to survive, with pass/fail
rules written down in advance. The reasoning: a researcher who decides what counts as
failure \emph{after} seeing results will always find a reason to keep the result.

\begin{itemize}[leftmargin=1.8em,itemsep=1pt]
  \item \textbf{The placebo test.} Re-run the entire design pretending the ban happened in
        \textbf{December 2019} - a date when nothing happened - using only pre-ban data.
        A trustworthy design should find \emph{nothing}. This is the one test where I am
        rooting for a null: finding an ``effect'' at a fake date means the machinery
        manufactures effects.
  \item \textbf{The pre-trend test.} An \textbf{event study} estimates the banned-vs-control
        gap month by month around the ban (the months before are called \emph{leads}, after
        are \emph{lags}). If the groups were already diverging \emph{before} the ban, the
        parallel-trends assumption is broken. Rule: if the leads are jointly significant,
        the DiD is dead regardless of how pretty its headline is.
\end{itemize}

\subsection{Checkpoint: the design fails its own tests}

First proper run on the district panel: the DiD said banned commodities' volatility was
about \textbf{16\% lower} post-ban ($p=0.029$) - opposite in sign to the popular fear.
Then the battery:

\begin{checkpoint}{Both alarms fired}
The fake-2019 placebo returned $\mathbf{-13.6\%}$, $p=0.046$ - statistically significant,
at a date where nothing happened, and nearly the size of the ``real'' effect. The joint
pre-trend test rejected ($p=0.033$). By my own pre-registered rules, the design was dead.
\end{checkpoint}

The tempting move was to soften the rules or quietly drop the placebo. Instead I treated
the failure as information: something upstream was manufacturing effects at arbitrary
dates. That points at \emph{measurement}, not economics.

\subsection{Checkpoint: the calendar-grid artifact}

\begin{checkpoint}{The data error the placebo exposed}
The mandi price series were \textbf{seven-day calendar grids}: across hundreds of
districts, some market reports on nearly every calendar day, and closed days simply carry
Friday's price forward - in one series, 28\% of rows were weekend dates. So my ``daily
returns'' included Fri$\to$Sat$\to$Sun steps computed on carried, unchanged prices (fake
zero returns and weekend autocorrelation), while my volatility formula annualised by
$\sqrt{252}$ \emph{trading} days - a units mismatch between the sampling grid and the
scaling constant. Every volatility number in the panel was contaminated.
\end{checkpoint}

The fix is the one-line weekday filter shown in Part 3, followed by a full rebuild of the
panel and every estimate. The proof it was the true cause: after the fix, \textbf{the
placebo dropped below significance and the pre-trends went parallel} - the design
survived its own rules - while the headline moved only moderately. (Full honesty: my
internal audit still graded that placebo \emph{marginal} - insignificant, but recovering
roughly half the headline's size - a residual worry that only the better comparison
group of Part 4.8 fully cleared.) And a side-result died honourably: an early
GARCH fit had shown chana's volatility persistence collapsing from 0.99 to 0.78 post-ban -
dramatic, quotable, and entirely an artifact of those fake weekend zero-returns. On clean
data it is roughly flat. I retired it.

\begin{lesson}
The falsification battery did exactly the job it was designed for - it caught a bug I did
not know I had, in my own pipeline, before it could reach a headline. Pre-registering the
tests is what made the failure undeniable.
\end{lesson}

\subsection{Checkpoint: honest uncertainty with few treated units}

On clean data (and with paddy excluded for MSP censoring) the DiD read $\mathbf{-20.7\%}$.
The naive clustered $p$-value said 0.008 - but that number leans on large-sample theory
that needs many clusters, and at this stage I had ten commodity clusters, only five of
them treated.
Two corrections give honest uncertainty:

\begin{itemize}[leftmargin=1.8em,itemsep=1pt]
  \item Refer the cluster-robust $t$-statistic to a $t$ distribution with $G-1$ degrees of
        freedom (with $G$ = number of clusters) - the small-sample correction. Result:
        $p \approx 0.026$.
  \item The \textbf{wild-cluster bootstrap}: impose the null (no effect), then rebuild the
        outcome thousands of times by randomly flipping the sign of each \emph{whole
        commodity's} residuals ($\pm1$ with equal probability - Rademacher weights),
        re-estimating each time (999 draws; $p=(1+\#\{|t^*|\ge|t|\})/(B+1)$). The spread
        of those fake $t$-statistics is the honest yardstick. Result: $p \approx 0.046$.
\end{itemize}

So: borderline significant, and reported as such - the 0.008 appears only as a cautionary
example of fake precision.

\subsection{Checkpoint: the second design layer - synthetic twins}

Because the ban's timing was itself a reaction to high prices (governments ban when
markets run hot), a lingering worry is \textbf{mean reversion}: hot markets cool on their
own. Two further estimators attack this by building an explicit counterfactual:

\begin{itemize}[leftmargin=1.8em,itemsep=1pt]
  \item \textbf{Synthetic control (SCM):} for each banned commodity, find non-negative
        weights (summing to one) over the comparison commodities so the weighted blend -
        the ``synthetic twin'' - tracks the banned commodity's \emph{pre-ban} volatility
        path as closely as possible. The post-ban gap between commodity and twin is the
        effect. Significance comes from \textbf{in-space placebos}: pretend each comparison
        commodity was treated, fit its twin, and ask whether the real treated gap stands
        out from that placebo distribution.
  \item \textbf{Synthetic DiD:} a hybrid that also reweights \emph{time periods} and adds
        an intercept, so the twin need only be parallel, not identical - it relaxes SCM's
        harshest requirement.
\end{itemize}

An internal audit of this stage caught two inference errors I had made and am glad to own:
a pooled synthetic-DiD ``$z=-7.7$'' that turned out to measure donor-swap stability rather
than genuine sampling uncertainty (withdrawn - the honest pooled inference is borderline),
and a placebo comparison at the district level that compared a mean against single noisy
units (fixed to like-for-like). Both corrections made the results \emph{weaker} and the
paper \emph{stronger}.

\subsection{Checkpoint: the comparison group changes the answer}

\begin{checkpoint}{The counterfactual revision - the most important finding}
My original comparison commodities - castor, guar, cotton, jeera, turmeric - were chosen
because their futures kept trading. But they are industrial, fibre, and spice crops; the
banned commodities are food staples. When I extended the data collection to non-banned
\emph{food} crops (barley, maize, jowar, bajra) and added them as comparisons, the
estimate \textbf{halved: $-20.7\%$ to $-9.8\%$} - and lost statistical significance
(few-cluster $p\approx0.145$; bootstrap $0.153$). The food controls' volatility had also
declined over 2022-25; part of my earlier ``effect'' was simply the difference between
food and non-food markets.
\end{checkpoint}

Under the better comparison group the falsification battery got \emph{cleaner} (placebo
$-6.5\%$, $p\approx0.34$; pre-trends $p\approx0.45$), and the commodity detail became the
real story: \textbf{chana} is the one commodity with a strong, individually significant
decline under every estimator (synthetic DiD $-38.7\%$, $z=-2.05$); \textbf{wheat}
collapsed to approximately zero ($+0.4\%$) - its apparent decline had been the export
bans, stock limits, and procurement operations that hit wheat simultaneously, exactly what
the policy ledger existed to flag. Leave-one-commodity-out confirms the pattern: drop
wheat and the effect \emph{strengthens} to $-15.2\%$ ($p\approx0.003$); drop chana and it
thins to $-7.4\%$.

\begin{lesson}
The estimate answers to the quality of its counterfactual. I treat the halving not as an
embarrassment but as the project's most informative single result - and I would rather
carry a modest number I can defend than a large one I cannot.
\end{lesson}

\subsection{The cost side, and the verdict}

Three companion analyses complete the picture:

\begin{itemize}[leftmargin=1.8em,itemsep=1pt]
  \item \textbf{Who used these markets?} Parsing SEBI's participant tables: in the year
        before the ban, registered hedgers were about \textbf{1.4\%} of agricultural
        turnover (42\% proprietary, 57\% client). ``Farmers lost their hedging tool'' is
        not where the cost lies - direct hedging barely existed.
  \item \textbf{Market integration.} \textbf{Cointegration} - a statistical test of
        whether two price series are tied together over the long run, here applied across
        mandis - shows banned commodities' markets became measurably \emph{less}
        integrated post-ban (share of cointegrated market pairs fell 0.22 more than for
        comparisons, one-sided $p\approx0.07$), with slower price-gap correction. The
        comparisons \emph{improved} over the same years, so this is if anything an
        understatement.
  \item \textbf{Volatility dynamics.} A structural-break scan (with the
        heteroskedasticity-robust detector - the textbook one over-fires on fat-tailed
        returns) finds \textbf{no volatility-regime break near the ban date} in any banned
        commodity.
\end{itemize}

\begin{checkpoint}{The final verdict, honestly stated}
The feared volatility \emph{rise} is refuted. The measured decline is about $-10\%$,
statistically insignificant under honest few-cluster inference, and concentrated in chana;
wheat's decline is confounded and not attributable to the futures ban. The demonstrable
cost of the ban is informational - reduced spot-market efficiency and integration - not
lost hedging. Suggestive, not conclusive; and every number above traces to a file in the
project.
\end{checkpoint}

% =====================================================================
\section{The checkpoint map (one-page summary of the journey)}

\begin{center}\small
\begin{tabular}{p{0.5cm} p{4.6cm} p{9.2cm}}
\toprule
\# & Checkpoint & What happened / what it taught \\
\midrule
1 & Question + rules fixed & Evaluation question set; raw-immutability, logging, pre-registration, falsification rules locked before data. \\
2 & Route to mandi prices & Official portal unusable in bulk; documented API route found; provenance dual-cited. \\
3 & Rate-limit war & Naive puller burned the quota (17 min, 0 rows); rebuilt to read rate-limit headers and sleep to reset. \\
4 & Resumable architecture & One self-describing JSON per work unit; skip-if-exists; backoff on socket errors; 5.74M rows over $\sim$10 h. \\
5 & 504 divide-and-conquer & Oversized state requests split by district batches, then halved, then date-halved for one dense district. \\
6 & Identifier + robot traps & API ids $\ne$ portal ids (verified by name); exchange 403s beaten via a real browser on its own endpoint. \\
7 & Ghost prices & Post-ban settlement marks with zero volume $\Rightarrow$ hard truncation rule vindicated. \\
8 & Guar forensics & Series labelled guar correlated 0.04 with guar futures - gum contamination; correct series (0.99) swapped in. \\
9 & Cleaning pipeline & Log returns, rolling $\sqrt{252}$ volatility, median aggregation, MSP censoring (paddy out, wheat flagged). \\
10 & First estimates & District DiD $-16.4\%$ ($p=.029$) - opposite of the feared rise. \\
11 & \textbf{Design fails own tests} & Placebo $-13.6\%$ ($p=.046$) at a fake date; pre-trends reject ($p=.033$). Dead by my own rules. \\
12 & \textbf{Calendar-grid artifact} & Weekend rows + $\sqrt{252}$ mismatch found and fixed; battery flips to passing; a flashy GARCH result evaporates. \\
13 & Honest inference & $-20.7\%$ with few-cluster $p\approx.026$, bootstrap $.046$; naive $.008$ demoted to a cautionary tale. \\
14 & Synthetic twins + audit & SCM/SDID corroborate sign; two of my own inference errors caught and withdrawn. \\
15 & \textbf{Counterfactual revision} & Food-crop donors halve the estimate to $-9.8\%$, insignificant; chana concentrated; wheat $\approx$ 0 (confounded). \\
16 & Cost side & 1.4\% hedger share; integration loss ($-0.22$, $p\approx.07$); no volatility-regime break at the ban. \\
17 & Final verdict & Rise refuted; modest, chana-specific calming; cost = price discovery. Everything documented and reproducible. \\
\bottomrule
\end{tabular}
\end{center}

\end{document}
